\section{Token-level EEG 视觉特征增强}

在 A2 方案中，EEG 融合是在 pooled embedding 这一层完成的，融合之后得到一个 512 维的向量，图像的空間结构信息在这个过程中丢掉了。为了让 EEG 信息能够进入视觉语言模型原本的 token 处理流程，我们在本章提出 VTF 方法，直接在 CLIP ViT 的 visual patch token 这一层把 EEG 信息加进去。

\subsection{为什么需要 token 级融合}

Pooled embedding fusion 的做法是把整张图压缩成一个向量再和 EEG 融合，这种做法做分类没什么问题，但到了生成式任务就有明显的短板：语言模型需要的输入是一个 token 序列，而不是一个单独的向量；token 序列里保留了图像不同空间位置的特征，这样才能对图中的多个对象做细致描述。另外，EEG 对不同区域的响应可能是不一样的——比如看一张同时有人物和物体的复杂图片时，各脑区的响应模式会有区别。Token-level fusion 的思路是用注意力机制让不同的 visual patch token 各自选择性地关注不同的 EEG token，从而实现更细粒度的跨模态交互。

\subsection{从 CLIP ViT 提取视觉 token}

CLIP ViT-B/32 的处理流程是先把 $224\times224$ 的图像切成 $7\times7$ 个 patch（每个 patch 是 $32\times32$），经过 ViT encoder 之后输出 50 个视觉 token（包括 1 个 CLS token 和 49 个 patch token），每个 token 的维度是 512，即 $F_{\mathrm{vis}} \in \mathbb{R}^{B\times50\times512}$。其中 CLS token 负责汇总全局语义信息，patch token 负责表示每个局部区域的特征。

\subsection{构造 EEG token}

A2 encoder 输出的 pooled EEG embedding 是一个 512 维的向量。为了让它能和视觉 token 做 token 级别的交互，我们设计了一个轻量的 EEG tokenizer：用 4 个并行的可学习线性投影器把 pooled embedding 映射成 4 个 EEG token，得到 $F_{\mathrm{eeg}} \in \mathbb{R}^{B\times4\times512}$。之所以选 4 个 token，是我们在效率和信息丰富度之间做的权衡——token 太少的话交互多样性不够，太多又可能把视觉表示压得太厉害。

\subsection{用交叉注意力做视觉到 EEG 的融合}

VTF 的核心是一个 visual-to-EEG 的交叉注意力模块，整体结构如图~\ref{fig:vtf} 所示。我们把视觉 token 当作 Query，EEG token 当作 Key 和 Value：

\begin{equation}
  F_{\mathrm{enhanced}} = F_{\mathrm{vis}} + \beta \cdot \mathrm{CrossAttention}(Q=F_{\mathrm{vis}}, K=F_{\mathrm{eeg}}, V=F_{\mathrm{eeg}})
\end{equation}

\placeholderfigure{VTF visual-to-EEG cross-attention 结构图}{vtf-cross-attention-placeholder}{4.2cm}

这里面的 CrossAttention 用的是标准的多头交叉注意力机制，头数 $h=8$。$\beta$ 是一个置信度感知的门控系数，它的计算方式如下：

\begin{equation}
  \beta = \gamma \cdot \sigma\left(W_\beta [\mathrm{pool}(F_{\mathrm{vis}}); \mathrm{pool}(F_{\mathrm{eeg}})] + b_\beta\right)
\end{equation}

$\gamma$ 是一个可训练的缩放因子，初始值设为 0.1，这样在训练初期 EEG 的注入会被压制住，避免出现梯度不稳定的问题。这个门控机制的作用是根据当前视觉信号的质量来决定多大程度上引入 EEG 信息。融合后的 $F_{\mathrm{enhanced}}$ 形状仍然是 $[B, 50, 512]$，任何下游模型只要能接收 CLIP ViT 的 visual tokens 就可以直接用。

\subsection{用分类实验做验证}

为了确认 VTF 确实有效果，我们先拿分类任务来做评估：把 $F_{\mathrm{enhanced}}$ 做平均池化，然后用 CLIP 的文本原型做零样本分类。实验结果显示，在多种退化条件下，real EEG 的效果明显好于只用视觉信号（vision-only）以及用 shuffled 或 random EEG 做对照的情况，这说明 token-level 的 EEG fusion 确实是可行的。不过话说回来，VTF 的整体分类准确率还是明显不如 A2 的 pooled fusion，原因在于 A2 用了更深的 gated fusion 网络，跨模态交互也更充分。VTF 的真正价值在于它生成的 enhanced vision tokens 可以直接喂给生成式 VLM，这是下一章要做的事情。
